% =====================================================================
%  rosei_master.tex
%  Master rederivation of Rosei's interband epsilon_2 for gold from
%  Fermi's Golden Rule in k-space, written throughout in Rosei's own
%  notation and symbols (GRW 1975), so that his ten numbered equations
%  pop out of one derivation.  Where the printed equations do not come
%  out -- (4), (6), (8) -- the failure is stated and fixed.  Includes
%  the photonic density of states, the emission counterpart, and the
%  k_parallel quadrature scheme used by the numerics.
%
%  Consolidates and supersedes (see derivations/README.md):
%    A1, A8, A8.L, "Absorption Integral (Using Rosei's Notations)",
%    "Interband Absorption at X and L (Rosei's Notation)",
%    "Equivalence of My Absorption Integral and Rosei's",
%    "Effective Mass Extraction", rosei_model_formulas.md,
%    interband_derivation{,_v2}.tex, rosei_equivalence_derivation.tex;
%    notation and fixes follow "Rebuilding Rosei's Interband Absorption
%    in His Own Notation & Suggesting Possible Fixes.md".
% =====================================================================
\documentclass[11pt,a4paper]{article}

\usepackage[hmargin=2.0cm,vmargin=2.4cm]{geometry}
\usepackage{amsmath,amssymb,mathtools,bm}
\usepackage{booktabs}
\usepackage{paracol}
\usepackage{needspace}
\usepackage[dvipsnames,table]{xcolor}
\usepackage{microtype}
\usepackage[colorlinks=true,linkcolor=MidnightBlue,citecolor=ForestGreen,urlcolor=MidnightBlue]{hyperref}

\numberwithin{equation}{section}

% ------------------------------ macros -------------------------------
\newcommand{\hw}{\hbar\omega}
\newcommand{\Eg}{\mathcal{E}_g}
\newcommand{\EgX}{\mathcal{E}_g^{X}}
\newcommand{\EgL}{\mathcal{E}_g^{L}}
\newcommand{\Elu}{\mathcal{E}_{lu}}
\newcommand{\cD}{\mathcal{D}}
\newcommand{\cF}{\mathcal{F}}
\newcommand{\cJ}{\mathcal{J}}
\newcommand{\Ab}{\bar{A}}
\newcommand{\Bb}{\bar{B}}
\newcommand{\kperp}{k_{\perp}}
\newcommand{\kpar}{k_{\parallel}}
\newcommand{\pul}{\mathbf{p}_{ul}}
\newcommand{\eps}{\varepsilon}
\newcommand{\dd}{\mathrm{d}}
\newcommand{\Dlu}{D_{l\to u}}
\newcommand{\wXu}{\hbar\omega_{X_6^-}}
\newcommand{\wXl}{\hbar\omega_{X_7^+}}
\newcommand{\wLu}{\hbar\omega_{L_6^-}}
\newcommand{\wLl}{\hbar\omega_{L^+}}
\newcommand{\wuz}{\hbar\omega_{u0}}
\newcommand{\wlz}{\hbar\omega_{l0}}
\newcommand{\fFD}{f}
\newcommand{\kB}{k_B}
\newcommand{\rme}{\mathrm{e}}

% step titles inside the two-column X/L comparison
% (vspace* so the gap survives at the top of a synchronized row)
\newcommand{\colstep}[1]{\par\vspace*{10pt}\noindent\textbf{#1}\par\nobreak\vspace*{3pt}\noindent\ignorespaces}

% shading for the core rows of the route map
\colorlet{corerow}{MidnightBlue!8}

\newtheorem{remark}{Remark}[section]

\title{\bfseries Interband Absorption in Gold from Fermi's Golden Rule:\\
a first-principles rederivation of Rosei's $\eps_2$ at the $X$ and $L$
critical points, in Rosei's own notation}
\author{Ben Shpigel\\ \small Department of Physics, Ben-Gurion University of the Negev}
\date{July 2026}

\begin{document}
\maketitle

\begin{abstract}
\noindent
This document rederives, from first principles and with every prefactor
tracked, the Rosei-model interband ${\varepsilon}_2({\hbar\omega},T)$ of gold at the $X$
and $L$ critical points [Guerrisi, Rosei \& Winsemius, Phys.\ Rev.\ B
\textbf{12}, 557 (1975); hereafter simply ``Rosei''], written
throughout in Rosei's own notation, so that the paper's ten numbered
equations emerge explicitly from one derivation.  The route runs from
Fermi's Golden Rule (FGR) for Bloch states in $\mathbf{k}$-space,
through the reduction of the transition sum to a one-dimensional
energy integral by the change of variables $(u,v)=({k_{\perp}}^2,{k_{\parallel}}^2)$,
to the resolution of the double $\delta$-constraint by a single
$2\times2$ Jacobian.  Seven of the ten equations are recreated
verbatim; the printed Eqs.~(4), (6) and (8) do not come out as
written, and for each the failure is stated in one line and the
corrected form derived (Sec.~\ref{sec:grw}).  The photonic density of
states is derived by field quantization, the spontaneous-emission
counterpart is assembled, and the pair is shown to close exactly into
the detailed-balance relation between emission and absorption (van
Roosbroeck--Shockley/Kirchhoff) --- the global consistency check of
all prefactors.  A final section gives the ${k_{\parallel}}$ quadrature scheme
that evaluates the model without touching its edge singularities: the
groundwork for the numerical implementation.
\end{abstract}
\tableofcontents

% =====================================================================
\section{Scope \& sources}\label{sec:scope}
% =====================================================================

Everything below is self-contained: it is built from Fermi's Golden
Rule and four primary sources.
\begin{itemize}
\item \textbf{Rosei 1975} (Guerrisi, Rosei \& Winsemius; ``Rosei'' from here on) --- the paper recreated.  Only the $X$ point is actually derived there, and the printed Eqs.~(4), (6), (8) have defects --- fixed in Sec.~\ref{sec:grw}.
\item \textbf{Rosei 1974} --- the EDJDOS recipe and the $L$-point machinery (``will not be repeated here''; it is repeated here).
\item \textbf{Christensen \& Seraphin 1971} --- the band structure: levels and masses.  Rosei's own source as well (his Ref.~4).
\item \textbf{Johnson \& Christy 1972} --- the optical data the model is fitted to.
\end{itemize}
Notation is Rosei's: bands are $u$ (upper) and $l$ (lower), all levels
and masses are positive numbers with every sign written explicitly,
and an equation tagged \textbf{($n$)} is Rosei's Eq.~($n$) recreated
--- primes, as in ($6'$), mark corrected forms.  The rest of the
conventions live in App.~\ref{sec:conv}.

\medskip
\noindent Route map --- shaded rows are the heart of the document, the
Rosei recreation; the others are scaffolding that keeps it
self-contained:

\begin{center}
\small
\renewcommand{\arraystretch}{1.3}
\begin{tabular}{@{}lll@{}}
\toprule
where & what happens & leans on \\
\midrule
Sec.~\ref{sec:fgr} & golden-rule machinery, every constant tracked & --- \\
\rowcolor{corerow}
Sec.~\ref{sec:reduction} & $\mathbf k$-space reduction $\to$ Eqs.~($4'$), (5), (7) & Rosei 1975/1974 \\
\rowcolor{corerow}
Sec.~\ref{sec:XL} & $X$ and $L$, side by side $\to$ Eqs.~($6'$), ($8'$) & Rosei 1975/1974 \\
\rowcolor{corerow}
Sec.~\ref{sec:eps2} & ${\varepsilon}_2$ assembled $\to$ (9)--(11); scorecard; corrections; Drude & C\&S; J\&C \\
Sec.~\ref{sec:photonic} & photonic DOS, emission, detailed balance & Novotny--Hecht \\
Sec.~\ref{sec:numerics} & ${k_{\parallel}}$ quadrature for the implementation & --- \\
\bottomrule
\end{tabular}\end{center}

% =====================================================================
\section{Fermi's Golden Rule in $\mathbf{k}$-space}\label{sec:fgr}
% =====================================================================

\subsection{Minimal coupling}

With the radiation field described by a vector potential $\mathbf{A}$
in the Coulomb gauge ($\nabla\!\cdot\!\mathbf{A}=0$), the one-electron
Hamiltonian is
\begin{equation}
\hat H=\frac{[\hat{\mathbf p}+\tfrac{e}{c}\mathbf A(\hat{\mathbf r},t)]^2}{2m}
       +V(\hat{\mathbf r})
      =\underbrace{\frac{\hat p^2}{2m}+V}_{\hat H_0}
       +\underbrace{\frac{e}{2mc}\bigl(\mathbf A\!\cdot\!\hat{\mathbf p}
        +\hat{\mathbf p}\!\cdot\!\mathbf A\bigr)}_{\hat H_{\rm int}}
       +\underbrace{\frac{e^2A^2}{2mc^2}}_{\text{dropped}} .
\end{equation}
In the Coulomb gauge $\hat{\mathbf p}\!\cdot\!\mathbf A
=\mathbf A\!\cdot\!\hat{\mathbf p}$, so
\begin{equation}
\boxed{\;\hat H_{\rm int}=\frac{e}{mc}\,\mathbf A\!\cdot\!\hat{\mathbf p}\;}
\label{eq:Hint}
\end{equation}
(the $\tfrac12$ of the symmetrized form is consumed by the identity of
the two terms).  The $A^2$ term does not connect different bands to
first order and is dropped.

\subsection{Quantized field: absorption, stimulated and spontaneous emission}
\label{sec:quantized}

Expanding the field in a quantization volume $V$,
\begin{equation}
\mathbf A(\mathbf r)=\sum_{\mathbf q,\lambda}
\sqrt{\frac{2\pi\hbar c^2}{\omega_q V}}\;\hat{\bm{\varepsilon}}_\lambda
\Bigl(\hat a_{\mathbf q\lambda}\,{\mathrm{e}}^{i\mathbf q\cdot\mathbf r}
     +\hat a^\dagger_{\mathbf q\lambda}\,{\mathrm{e}}^{-i\mathbf q\cdot\mathbf r}\Bigr),
\label{eq:Aquant}
\end{equation}
the matrix elements of $\hat H_{\rm int}$ between joint
electron--photon states carry the boson factors
\begin{equation}
\bigl|\langle n_q-1|\hat a|n_q\rangle\bigr|^2=n_q
\quad\text{(absorption)},\qquad
\bigl|\langle n_q+1|\hat a^\dagger|n_q\rangle\bigr|^2=n_q+1
\quad\text{(emission)} .
\label{eq:bosefactors}
\end{equation}
The ``$+1$'' is the spontaneous channel: it survives at $n_q=0$ and is
invisible to a purely classical $\mathbf A$.  The semiclassical golden
rule gives absorption and \emph{stimulated} emission only; spontaneous
emission requires \eqref{eq:bosefactors}, and its mode sum produces the
photonic density of states in Sec.~\ref{sec:photonic}.

\subsection{Bloch states and vertical transitions}

For Bloch states $\psi_{n\mathbf k}={\mathrm{e}}^{i\mathbf k\cdot\mathbf r}
u_{n\mathbf k}(\mathbf r)/\sqrt V$, the matrix element of
${\mathrm{e}}^{\pm i\mathbf q\cdot\mathbf r}\,\hat{\bm{\varepsilon}}\!\cdot\!\hat{\mathbf p}$
enforces $\mathbf k'=\mathbf k\pm\mathbf q$.  At optical frequencies
$q\sim\omega/c\sim10^{-3}\,$\AA$^{-1}$ is negligible on the zone scale,
so transitions are vertical ($\mathbf k'=\mathbf k$) and
\begin{equation}
M_{ul}(\mathbf k)\equiv
\langle u\,\mathbf k|\,\hat{\bm{\varepsilon}}\!\cdot\!\hat{\mathbf p}\,|l\,\mathbf k\rangle,
\qquad
\bigl|\langle f|\hat H_{\rm int}|i\rangle\bigr|^2
=\Bigl(\frac{e}{mc}\Bigr)^{2}\frac{2\pi\hbar c^2}{\omega V}\,
 |M_{ul}(\mathbf k)|^{2}\times\{n_q\text{ or }n_q+1\}.
\label{eq:MEsq}
\end{equation}

\subsection{Golden-rule rate, ${\varepsilon}_2$, and the origin of the $1/\omega^2$}
\label{sec:eps2def}

FGR for one photon mode, then summing over $\mathbf k$ (factor 2 for
spin, $\sum_{\mathbf k}\to V\!\int\!{\mathrm{d}}^3k/(2\pi)^3$) with Pauli
blocking, gives the photon absorption rate.  Equating the absorbed
power to the electromagnetic dissipation rate
$P_{\rm abs}=\tfrac{1}{8\pi}\,\omega\,{\varepsilon}_2(\omega)\,|\mathbf E_0|^2 V$
of a classical field $\mathbf E=\mathbf E_0\cos\omega t$ yields
\begin{equation}
\boxed{\;
{\varepsilon}_2(\omega)=\frac{4\pi^2 e^2}{m^2\omega^2}\,
\frac{2}{(2\pi)^3}\int_{\rm BZ}\!{\mathrm{d}}^3k\;
\bigl|M_{ul}(\mathbf k)\bigr|^{2}\,
\delta\!\bigl(\hbar\omega_u(\mathbf k)-\hbar\omega_l(\mathbf k)-{\hbar\omega}\bigr)\,
\bigl[1-{f}\bigl(\hbar\omega_u(\mathbf k)\bigr)\bigr]\; }
\label{eq:eps2exact}
\end{equation}
per pair of bands (SI: replace $4\pi e^2\to e^2/{\varepsilon}_0$).  The
occupation factor is Rosei's, discussed in Sec.~\ref{sec:occupation}.

Where does the $1/\omega^2$ up front come from?  It is fixed before
any band structure enters --- and it is worth watching, because
emission will carry a \emph{different} power
(Sec.~\ref{sec:photonic}).  Unpack \eqref{eq:eps2exact} into its three
$\omega$-carrying pieces, marking the net power each one contributes:
\begin{equation}
{\varepsilon}_2=\frac{8\pi\,P_{\rm abs}/V}
{\underbrace{\textstyle\omega}_{-1}\;E_0^{2}},
\qquad
P_{\rm abs}=\underbrace{\textstyle{\hbar\omega}}_{+1}\times\ \text{rate},
\qquad
\text{rate}\;\propto\;
\Bigl(\frac{eA_0}{mc}\Bigr)^{\!2}
=\frac{e^{2}E_0^{2}}{m^{2}}\,
\underbrace{\frac{1}{\omega^{2}}}_{-2}\ \ (A_0=cE_0/\omega).
\label{eq:wcount}
\end{equation}
The $+1$ (one quantum per event) cancels the $-1$ (the $\omega$ in the
definition of dissipation); the $-2$ of the vertex is all that
survives.  It is the price of writing the
$\mathbf A\!\cdot\!\mathbf p$ coupling in terms of the measured field
--- pure kinematics, no density of states involved.  (Length-gauge
cross-check: $|{\mathbf{p}_{ul}}|^2=m^2\omega_{ul}^2\,|\mathbf r_{ul}|^2$ hides the
same factor inside the dipole element, so on shell the two forms agree
identically.)

Near a critical point the interband matrix element varies slowly;
following Rosei we freeze it and average over polarization, in his
gradient convention:
\begin{equation}
|P|^2\equiv\bigl|\langle u|\nabla|l\rangle\bigr|^2
\ \ (\text{length}^{-2}),
\qquad
\overline{|M_{ul}|^2}
=\overline{\bigl|\hat{\bm{\varepsilon}}\!\cdot\!{\mathbf{p}_{ul}}\bigr|^2}
=\frac{\hbar^2|P|^2}{3} .
\label{eq:Pdef}
\end{equation}
Per critical-point family this turns \eqref{eq:eps2exact} into the
working form
\begin{equation}
{\varepsilon}_2({\hbar\omega},T)=\frac{4\pi^2e^2\hbar^2}{3m^2\omega^2}\cdot
\frac{2}{(2\pi)^3}\,N\,|P|^2
\int_{\text{half-nbhd}}\!\!{\mathrm{d}}^3k\;
\delta\bigl({\mathcal{E}_{lu}}(\mathbf k)-{\hbar\omega}\bigr)\,
\bigl[1-{f}(E)\bigr]\Bigr|_{E=\hbar\omega_u(\mathbf k)},
\label{eq:star}
\end{equation}
with the interband energy ${\mathcal{E}_{lu}}\equiv\hbar\omega_u-\hbar\omega_l$
(Sec.~\ref{sec:bands} connects it to Rosei's $\Omega_{lu}$) and $N$
the half-neighborhood count of App.~\ref{sec:conv}.

\subsection{Occupation factor: Rosei's $[1-f]$}\label{sec:occupation}

The exact net weight of \eqref{eq:eps2exact} --- absorption minus
stimulated emission --- is the difference
${f}(E-{\hbar\omega})-{f}(E)$ read at the final-state energy
$E=\hbar\omega_u(\mathbf k)$.  For the noble-metal $d$ bands the
initial state sits at $E-{\hbar\omega}\lesssim-1.5\,$eV${}\ll-{k_B} T$, so
${f}(E-{\hbar\omega})=1$ to machine precision and the difference \emph{is}
Rosei's factor,
\begin{equation}
{f}(E-{\hbar\omega})-{f}(E)\;\longrightarrow\;\bigl[1-{f}(E,T)\bigr],
\label{eq:occup}
\end{equation}
which is used throughout this document, exactly as Rosei uses it.
(The \emph{emission} weight is a different object --- a product, not a
difference --- and is derived in Sec.~\ref{sec:photonic}.)

% =====================================================================
\section{Reduction of the $\mathbf{k}$-integral: the common machinery}
\label{sec:reduction}
% =====================================================================

\subsection{The Rosei two-band model}\label{sec:bands}

Both critical points are axially symmetric parabolic pairs.  The
parabolic dispersion relations are the primary sources': at $X$ they
are Rosei's own Eqs.~(1)--(2); at $L$ they are the forms of the 1974
Ag analysis, which the 1975 paper adopts for gold; levels and masses
for both come from the relativistic band calculation of Christensen \&
Seraphin --- Rosei's source too.  Why label the bands by $E$ and
$E-{\hbar\omega}$?  Because a vertical transition at photon energy ${\hbar\omega}$
connects exactly those two energies, so the occupation factors can be
read off directly at the final energy $E$.  With the positive
levels/masses of App.~\ref{sec:conv}, generic levels ${\hbar\omega_{u0}}>0$ (empty,
above $E_F$) and ${\hbar\omega_{l0}}>0$ (occupied $d$ level below $E_F$), and the
\emph{topology switch} $s$,
\begin{equation}
\boxed{\;
\begin{aligned}
E=\hbar\omega_u(\mathbf k)&=+{\hbar\omega_{u0}}+A_u{k_{\perp}}^{2}+s\,B_u{k_{\parallel}}^{2},
&\qquad s_{X}&=-1\ \ \text{(saddle)},\\[2pt]
E-{\hbar\omega}=\hbar\omega_l(\mathbf k)&=-{\hbar\omega_{l0}}-A_l{k_{\perp}}^{2}-B_l{k_{\parallel}}^{2},
&\qquad s_{L}&=+1\ \ \text{(minimum)},
\end{aligned}\;}
\label{eq:bands}
\end{equation}
where $({\hbar\omega_{u0}},{\hbar\omega_{l0}})=({\hbar\omega_{X_6^-}},{\hbar\omega_{X_7^+}})$ at $X$ and $({\hbar\omega_{L_6^-}},{\hbar\omega_{L^+}})$ at $L$; the
lower band is a local maximum in both directions at both points.  The
interband energy is then
\begin{equation}
{\mathcal{E}_{lu}}(\mathbf k)
={\mathcal{E}_g}+{\bar{A}}\,{k_{\perp}}^{2}+{\bar{B}}\,{k_{\parallel}}^{2},
\qquad
{\mathcal{E}_g}={\hbar\omega_{u0}}+{\hbar\omega_{l0}},\quad
{\bar{A}}=A_u+A_l>0,\quad
\boxed{{\bar{B}}=B_l+s\,B_u}
\label{eq:Omega}
\end{equation}
--- the entire topological difference between $X$ and $L$ is carried
by $s$ through ${\bar{B}}$ (and the determinant ${\mathcal{D}}$ below).

A note on reach: \eqref{eq:bands} are local expansions, and the
constant matrix elements behind Eq.~(9) are claimed by Rosei only in a
region ${\sim}(\pi/10a)^{3}$ around each point --- yet no explicit
$k_{\max}$ will appear in any integration bound below.  That is
Rosei's own bookkeeping, not an omission: what polices locality
instead (the photon-energy range) is quantified after the window table
of Sec.~\ref{sec:eps2}.

A note on symbols: Rosei's own $\Omega_{lu}$ is not this energy but
the \emph{constraint} built from it,
\begin{equation}
\Omega_{lu}(\mathbf k)\equiv
\hbar\omega_u-\hbar\omega_l-{\hbar\omega}={\mathcal{E}_{lu}}(\mathbf k)-{\hbar\omega}=0,
\tag{3}
\label{eq:CEDS}
\end{equation}
whose zero set is the constant-energy-difference surface (CEDS) ---
the surface on which everything below happens.  We keep the two
apart: ${\mathcal{E}_{lu}}$ is a function on $\mathbf k$-space; Rosei's Eq.~(3),
$\Omega_{lu}=0$, picks out the surface.  The gold parameters are
collected in the tables of Sec.~\ref{sec:eps2}.

\subsection{Sorting transitions by final energy}
\label{sec:sort}

The occupation in \eqref{eq:star} depends on $\mathbf k$ only through
the final energy $E=\hbar\omega_u(\mathbf k)$ --- that is where
temperature enters --- so before any geometry we \emph{sort the
transitions by $E$}, by slipping the identity
$1=\int{\mathrm{d}} E\;\delta(\hbar\omega_u(\mathbf k)-E)$ into \eqref{eq:star}:
\begin{equation}
\begin{aligned}
\int_{{k_{\parallel}}>0}\!\!{\mathrm{d}}^3k\;\delta({\mathcal{E}_{lu}}-{\hbar\omega})\,[1-{f}]
&=\int_{E_{\min}}^{E_{\max}}\!\!{\mathrm{d}} E\;[1-{f}(E)]
\underbrace{\int_{{k_{\parallel}}>0}\!\!{\mathrm{d}}^3k\;
\delta\bigl(\hbar\omega_u-E\bigr)\,
\delta\bigl({\mathcal{E}_{lu}}-{\hbar\omega}\bigr)}_{(2\pi)^3\,{D_{l\to u}}(E,{\hbar\omega})}\\[2pt]
&\equiv(2\pi)^3\,{\mathcal{J}}_{l\to u}({\hbar\omega},T) .
\end{aligned}
\label{eq:sort}
\end{equation}
The limits $[E_{\min},E_{\max}]$ are simply where the double-$\delta$
has support; they are worked out per critical point in
Sec.~\ref{sec:XL}.  Read \eqref{eq:sort} from the outside in: the
outer object ${\mathcal{J}}_{l\to u}$ is already Rosei's thermally weighted JDOS
--- his Eq.~(7), waiting only for its integrand --- and that integrand,
the inner integral, is his central object: the
\emph{energy-distributed joint density of states} (EDJDOS),
transitions per unit crystal volume, per unit final energy, per unit
transition energy (one half-neighborhood, one spin).  Both $\delta$'s
are accounted for: energy conservation was already in \eqref{eq:star};
the second is the inserted identity, and it is what lets ${\mathcal{J}}$ weight
each $E$ by its own occupation.

\subsection{Change of variables $(u,v)=({k_{\perp}}^2,{k_{\parallel}}^2)$}\label{sec:linear}

This is the step where the present derivation simplifies Rosei's
geometric construction --- and the form of \eqref{eq:sort} itself
dictates it.  The two $\delta$-functions in \eqref{eq:sort} have
arguments $\hbar\omega_u(\mathbf k)-E$ and ${\mathcal{E}_{lu}}(\mathbf k)-{\hbar\omega}$, and
by \eqref{eq:bands} and \eqref{eq:Omega} both are \emph{affine in the
squared coordinates} ${k_{\perp}}^2$ and ${k_{\parallel}}^2$.  A $\delta$-function of
a linear argument collapses trivially; so choose the variables in
which the constraints \emph{are} linear.  On the quadrant $u,v\ge0$
with $(u,v)=({k_{\perp}}^2,{k_{\parallel}}^2)$, fixing the final energy $E$ and the
interband energy ${\mathcal{E}_{lu}}={\hbar\omega}$ reads
\begin{equation}
\begin{pmatrix} E-{\hbar\omega_{u0}}\\[2pt] {\hbar\omega}-{\mathcal{E}_g}\end{pmatrix}
=\underbrace{\begin{pmatrix} A_u & s\,B_u\\[2pt] {\bar{A}} & {\bar{B}}\end{pmatrix}}_{\textstyle M}
\begin{pmatrix} u\\[2pt] v\end{pmatrix},
\qquad
{\mathcal{D}}\equiv\det M=A_uB_l-s\,A_lB_u .
\label{eq:linsys}
\end{equation}
Stacking the two constraints \emph{is} the construction of the
transition matrix $M$; its determinant is the Jacobian that will
collapse the two $\delta$-functions below.  Inverting,
\begin{equation}
\boxed{\;
u=\frac{{\bar{B}}\,(E-{\hbar\omega_{u0}})-s\,B_u({\hbar\omega}-{\mathcal{E}_g})}{{\mathcal{D}}},
\qquad
v={k_{\parallel}}^{2}=\frac{A_u({\hbar\omega}-{\mathcal{E}_g})-{\bar{A}}\,(E-{\hbar\omega_{u0}})}{{\mathcal{D}}}
\;}
\label{eq:uv}
\end{equation}
and the physical window is exactly $\{u\ge0\}\cap\{v\ge0\}$.
The two point-specific determinants:
\begin{equation}
{\mathcal{D}}_X=A_uB_l+A_lB_u>0\ \ \text{always (either sign of ${\bar{B}}_X$)},
\qquad
{\mathcal{D}}_L=A_uB_l-A_lB_u\ \ (<0\ \text{for Au}) .
\label{eq:Ds}
\end{equation}

\subsection{The EDJDOS: general derivation for both $X$ and $L$}
\label{sec:kernel}

Now reduce the inner integral of \eqref{eq:sort} all the way --- this
is where ${\mathcal{D}}$ and ${\mathcal{F}}$ come from.
Axial symmetry gives ${\mathrm{d}}^3k=2\pi{k_{\perp}}{\mathrm{d}}{k_{\perp}}\,{\mathrm{d}}{k_{\parallel}}$; in the
squared variables ${k_{\perp}}{\mathrm{d}}{k_{\perp}}=\tfrac12{\mathrm{d}} u$,
${\mathrm{d}}{k_{\parallel}}={\mathrm{d}} v/(2\sqrt v)$:
\begin{equation}
\frac{1}{(2\pi)^3}\int_{{k_{\parallel}}>0}\!\!{\mathrm{d}}^3k
=\frac{1}{(2\pi)^3}\,\frac{\pi}{2}
\int_0^\infty\!\!{\mathrm{d}} u\int_0^\infty\!\frac{{\mathrm{d}} v}{\sqrt v}
=\frac{1}{16\pi^2}\int_0^\infty\!\!{\mathrm{d}} u\int_0^\infty\!\frac{{\mathrm{d}} v}{\sqrt v}.
\end{equation}
The two $\delta$'s are exactly the two rows of $M$ in
\eqref{eq:linsys}; linear $\delta$'s collapse onto the unique solution
\eqref{eq:uv} divided by $|\det M|=|{\mathcal{D}}|$, and $\sqrt{v_*}={k_{\parallel}}(E,{\hbar\omega})$:
\begin{equation}
{D_{l\to u}}(E,{\hbar\omega})=\frac{1}{16\pi^2\,|{\mathcal{D}}|}\,\frac{1}{{k_{\parallel}}(E,{\hbar\omega})} .
\label{eq:Draw}
\end{equation}
\textbf{This is the whole result of the reduction: it contains only
${\mathcal{D}}$ and ${k_{\parallel}}$ --- no ${\mathcal{F}}$ yet.}  Rosei's ${\mathcal{F}}_{l\to u}$ is not an
independent input but a repackaging of ${\mathcal{D}}$, manufactured here.
Expanding ${\mathcal{D}}_X$ in the masses,
\begin{equation}
{\mathcal{D}}_X=A_uB_l+A_lB_u
=\Bigl(\frac{\hbar^2}{2}\Bigr)^{2}
\frac{m_{l\perp}m_{u\parallel}+m_{l\parallel}m_{u\perp}}
{m_{l\perp}m_{l\parallel}m_{u\perp}m_{u\parallel}},
\end{equation}
and defining ${\mathcal{F}}_{l\to u}\equiv\hbar^2/(2\sqrt{|{\mathcal{D}}|})$ as pure
shorthand, the $\hbar^2/2$ cancels and Rosei's Eq.~(5) falls out ---
derived, not assumed --- together with its $L$ analogue:
\begin{equation}
{\mathcal{F}}_X=\left[\frac{m_{l\perp}m_{u\parallel}+m_{l\parallel}m_{u\perp}}
{m_{l\perp}m_{l\parallel}m_{u\perp}m_{u\parallel}}\right]^{-1/2},
\qquad
{\mathcal{F}}_L=\left[\frac{\bigl|m_{l\perp}m_{u\parallel}-m_{l\parallel}m_{u\perp}\bigr|}
{m_{l\perp}m_{l\parallel}m_{u\perp}m_{u\parallel}}\right]^{-1/2}
\tag{5}
\label{eq:Fmass}
\end{equation}
(the sum/difference structure is the $s=\mp1$ dichotomy of
\eqref{eq:Ds}).  Reading the shorthand backwards,
$|{\mathcal{D}}|=\hbar^4/4{\mathcal{F}}^2$, and \eqref{eq:Draw} becomes the EDJDOS in
Rosei's own symbols:
\begin{equation}
\boxed{\;
{D_{l\to u}}(E,{\hbar\omega})=\frac{1}{16\pi^2\,|{\mathcal{D}}|}\,\frac{1}{{k_{\parallel}}(E,{\hbar\omega})}
=\frac{{\mathcal{F}}_{l\to u}^{2}}{4\pi^2\hbar^{4}}\,\frac{1}{{k_{\parallel}}(E,{\hbar\omega})}
\;}
\tag{4$'$}
\label{eq:D4p}
\end{equation}
supported on $u,v\ge0$, with ${k_{\parallel}}$ from \eqref{eq:uv}.

\medskip
\noindent\textbf{Where the printed (4) falls.}  Rosei prints
${D_{l\to u}}=(8\pi^2\hbar^2)^{-1}{\mathcal{F}}/{k_{\parallel}}$.  With $|P|^2$ the squared
gradient matrix element of \eqref{eq:Pdef}, the pair \{printed (4),
printed (9)\} does not balance dimensionally, whereas ($4'$) ---
forced, with nothing left to choose, by the definition of ${D_{l\to u}}$ in
\eqref{eq:sort} --- closes the FGR chain exactly through to ${\varepsilon}_2$
(Sec.~\ref{sec:eps2}).  The two forms differ only by the constant
$2{\mathcal{F}}/\hbar^2$ per critical point; a constant cannot bend a line
shape, and it is absorbed into the fitted strength $S={\mathcal{F}}|P|^2$ of
Eq.~(10) --- which is why the slip is invisible in the published
numbers.  Details in Sec.~\ref{sec:grw}.

\subsection{The thermally weighted JDOS: Rosei's Eq.\ (7)}

With ${D_{l\to u}}$ in hand, the outer integral of \eqref{eq:sort} is,
verbatim as Rosei writes it,
\begin{equation}
\boxed{\;
{\mathcal{J}}_{l\to u}({\hbar\omega},T)=\int_{E_{\min}}^{E_{\max}}
{D_{l\to u}}(E,{\hbar\omega})\,\bigl[1-{f}(E,T)\bigr]\,{\mathrm{d}} E
\;}
\tag{7}
\label{eq:J7}
\end{equation}
with the windows $[E_{\min},E_{\max}]$ set by the geometry
($u,v\ge0$) and, where the geometry leaves an edge open, by the
occupation itself --- next section.

% =====================================================================
\section{The two critical points, side by side}\label{sec:XL}
% =====================================================================

Everything so far is common.  The topology forks the derivation
through two signs only: $s$ (shape of ${\bar{B}}$) and
$\mathrm{sign}({\mathcal{D}})$ (which window edge is singular).  The $X$ column
follows Rosei's own branch, ${\bar{B}}_X<0$ (see the mass fork,
Sec.~\ref{sec:fork}); his equations then pop out directly.

\Needspace{20\baselineskip}
\medskip
% \noindent\rule{\textwidth}{0.8pt}
\setlength{\columnsep}{22pt}
\columnratio{0.5}
\setlength{\columnseprule}{0.4pt}
\begin{paracol}{2}
\begin{center}\underline{{\Large\bfseries $X$ point\quad ($s=-1$)}}\end{center}
\switchcolumn
\begin{center}\underline{{\Large\bfseries $L$ point\quad ($s=+1$)}}\end{center}
\switchcolumn*
\colstep{Bands}
\begin{align*}
E&=+{\hbar\omega_{X_6^-}}+\frac{\hbar^2{k_{\perp}}^2}{2m_{u\perp}} -\frac{\hbar^2{k_{\parallel}}^2}{2m_{u\parallel}},\\ E-{\hbar\omega}&=-{\hbar\omega_{X_7^+}}-\frac{\hbar^2{k_{\perp}}^2}{2m_{l\perp}} -\frac{\hbar^2{k_{\parallel}}^2}{2m_{l\parallel}} .
\end{align*}
Rosei's Eqs.\ (1)--(2) verbatim: the $X_6^-$ ($sp$) band curves up in the face, down along $\Delta$ --- a saddle; $X_7^+$ ($d$) is a maximum.
\switchcolumn
\colstep{Bands}
\begin{align*}
E&=+{\hbar\omega_{L_6^-}}+\frac{\hbar^2{k_{\perp}}^2}{2m_{u\perp}} +\frac{\hbar^2{k_{\parallel}}^2}{2m_{u\parallel}},\\ E-{\hbar\omega}&=-{\hbar\omega_{L^+}}-\frac{\hbar^2{k_{\perp}}^2}{2m_{l\perp}} -\frac{\hbar^2{k_{\parallel}}^2}{2m_{l\parallel}} .
\end{align*}
Not derived in the 1975 paper (borrowed from Rosei 1974); rederived here with the same tools.  $L_6^-$ is a \emph{minimum} in both directions; $L^+$ ($d$ top) a maximum.  The only change from $X$ is the sign of the upper ${k_{\parallel}}^2$ term; that one flip cascades through everything below.
\switchcolumn*
\colstep{Transition surface (CEDS)}
\[{\mathcal{E}_{lu}}={\mathcal{E}_g^{X}}+{\bar{A}}_X{k_{\perp}}^2+{\bar{B}}_X{k_{\parallel}}^2,\quad {\bar{B}}_X=B_l-B_u<0\]
on Rosei's branch ($m_{u\parallel}<m_{l\parallel}$: the $sp$ band falls along $\Delta$ faster than the $d$ band).  The CEDS is an \emph{open} hyperboloid at every ${\hbar\omega}$: transitions exist on both sides of ${\mathcal{E}_g^{X}}$ --- a real sub-gap tail.
\switchcolumn
\colstep{Transition surface (CEDS)}
\[{\mathcal{E}_{lu}}={\mathcal{E}_g^{L}}+{\bar{A}}_L{k_{\perp}}^2+{\bar{B}}_L{k_{\parallel}}^2,\quad {\bar{B}}_L=B_l+B_u>0\]
unconditionally: a closed ellipsoid, existing only for ${\hbar\omega}\ge{\mathcal{E}_g^{L}}$. \emph{No sub-gap tail at $L$} --- the sharp-vs-tailed contrast of the two edges is baked in right here.
\switchcolumn*
\colstep{Determinant}
\[{\mathcal{D}}_X=A_uB_l+A_lB_u>0\ \ \text{(either branch)} .\]
\switchcolumn
\colstep{Determinant}
\[{\mathcal{D}}_L=A_uB_l-A_lB_u<0\ \ \text{for Au}\]
($m_{l\perp}m_{u\parallel}=0.084<m_{l\parallel}m_{u\perp}=0.247$).
\switchcolumn*
\colstep{Constrained ${k_{\parallel}
}$}
\[{k_{\parallel}}^{2} =\frac{2{\mathcal{F}}_X^{2}}{\hbar^{2}}\Bigl[ \frac{{\hbar\omega}-{\hbar\omega_{X_7^+}}-E}{m_{u\perp}} +\frac{{\hbar\omega_{X_6^-}}-E}{m_{l\perp}}\Bigr].\]
Rosei's Eq.\ (6), corrected --- called ($6'$) throughout; see Sec.~\ref{sec:grw}.
\switchcolumn
\colstep{Constrained ${k_{\parallel}
}$}
\[{k_{\parallel}}^{2} =\frac{2{\mathcal{F}}_L^{2}}{\hbar^{2}}\Bigl[ \frac{E-({\hbar\omega}-{\hbar\omega_{L^+}})}{m_{u\perp}} +\frac{E-{\hbar\omega_{L_6^-}}}{m_{l\perp}}\Bigr].\]
Same inversion, ${\mathcal{D}}_L<0$ reversing both inequalities.
\switchcolumn*
\colstep{Window}
$E$ decreases without bound along the CEDS (${\mathrm{d}} E\propto-{\mathcal{D}}_X{\mathrm{d}} v<0$): no geometric floor.  The integral is cut by occupation alone --- Rosei's practical floor $E_{\min}=-20{k_B} T$.  The ceiling is geometric and changes character at the gap:
\[E_{\max}={\hbar\omega_{X_6^-}}+ \begin{cases} \dfrac{A_u}{{\bar{A}}_X}({\hbar\omega}-{\mathcal{E}_g^{X}}), & {\hbar\omega}\ge{\mathcal{E}_g^{X}},\\[10pt] -\dfrac{B_u}{\vert{}{\bar{B}}_X\vert{}}({\mathcal{E}_g^{X}}-{\hbar\omega}), & {\hbar\omega}<{\mathcal{E}_g^{X}}, \end{cases}\]
the lower line being Rosei's Eq.\ (8), corrected --- ($8'$), Sec.~\ref{sec:grw}.
\switchcolumn
\colstep{Window}
Both edges geometric ($v\ge0$ now gives the \emph{lower} edge, $u\ge0$ the upper):
\begin{align*}
E_{\min}&={\hbar\omega_{L_6^-}}+\frac{A_u}{{\bar{A}}_L}({\hbar\omega}-{\mathcal{E}_g^{L}}),\\ E_{\max}&={\hbar\omega_{L_6^-}}+\frac{B_u}{{\bar{B}}_L}({\hbar\omega}-{\mathcal{E}_g^{L}}),
\end{align*}
slopes $0.745$ and $0.896$ for Au.  The whole window rides upward from ${\hbar\omega_{L_6^-}}=0.89$ eV (Remark~\ref{rem:neck}).
\switchcolumn*
\newpage
\colstep{Singular edge}
Above the gap the ceiling is the ${k_{\parallel}}\to0$ circle, where ${D_{l\to u}}\propto1/{k_{\parallel}}$ diverges:
\[\frac{1}{{k_{\parallel}}(E)}=\sqrt{\frac{{\mathcal{D}}_X}{{\bar{A}}_X}}\; \frac{1}{\sqrt{E_{\max}-E}} .\]
Below the gap the ceiling is the finite ${k_{\perp}}=0$ point: no divergence.
\switchcolumn
\colstep{Singular edge}
From \eqref{eq:uv} with ${\mathcal{D}}_L<0$,
\[{k_{\parallel}}^{2}=\frac{{\bar{A}}_L}{\vert{}{\mathcal{D}}_L\vert{}}\,\bigl(E-E_{\min}\bigr),\]
so ${D_{l\to u}}\propto1/{k_{\parallel}}$ diverges at the \emph{lower} limit:
\[\frac{1}{{k_{\parallel}}(E)}=\sqrt{\frac{\vert{}{\mathcal{D}}_L\vert{}}{{\bar{A}}_L}}\; \frac{1}{\sqrt{E-E_{\min}}} .\]
\switchcolumn*
\colstep{Onset}
At $T=0$ absorption starts when $E_{\max}$ crosses $E_F$:
\[{\hbar\omega}_{\rm on}={\mathcal{E}_g^{X}}-{\hbar\omega_{X_6^-}}\,\frac{m_{l\parallel}-m_{u\parallel}} {m_{l\parallel}}\simeq1.83\ \text{eV},\]
the piezomodulation threshold; the tail $1.83$--$1.94$ eV is purely geometric, and the onset is \emph{step-like} (Rosei's word), not divergent.
\switchcolumn
\colstep{Onset}
At ${\hbar\omega}={\mathcal{E}_g^{L}}$ exactly, with ${\mathcal{J}}_L\propto\sqrt{{\hbar\omega}-{\mathcal{E}_g^{L}}}$; the inverse-square-root weight sits on the onset edge itself, which is why the $L$ edge is sharp and strong.
\end{paracol}
\vspace{1em}
\bigskip
\noindent
\textbf{Punchline.}  The same machinery, forked only by two signs,
lands on two very different absorption edges.  At $L$ the $1/{k_{\parallel}}$
divergence of ${D_{l\to u}}$ sits on the \emph{onset} edge itself, so
absorption switches on sharply and strongly at ${\mathcal{E}_g^{L}}=2.45$ eV.  At
$X$ the divergence sits at the \emph{top} of the window, far from
onset; the onset itself is a soft, step-like rise near $1.83$ eV with
a geometric sub-gap tail. One interband edge that is really two ---
soft at $X$, sharp half an eV above it at $L$ --- and that is exactly
the ``splitting of the interband absorption edge'' of the 1975 title.

\begin{remark}[The $L$ level and the Fermi-surface neck]\label{rem:neck}
Rosei's qualitative discussion attributes the strength of the $L$
singularity to the onset CEDS being tangent to the Fermi surface along
the neck --- which requires $L_6^-$ \emph{below} $E_F$.  The fitted
scheme (${\mathcal{E}_g^{L}}=2.45$, ${\hbar\omega_{L^+}}=1.56\Rightarrow{\hbar\omega_{L_6^-}}=+0.89$ eV) puts the
whole $L$ window \emph{above} $E_F$, making the $L$ channel inert
through the Fermi factor.  Both statements cannot hold in one
parabolic model; the fitted scheme is what the code implements, and
the tension is recorded here rather than papered over.
\end{remark}
% ---------------------------------------------------------------------
\subsection{The $X$ mass fork: sign of ${\bar{B}}_X$}\label{sec:fork}
% ---------------------------------------------------------------------

One assignment at $X$ I am not sure of, flagged because it flips a
sign.  First, where the masses come from: they are not printed as
numbers anywhere --- they are read off the band diagram of C\&S
(Fig.~5), by fitting a parabola to the printed $E(\mathbf k)$ curves
around $X$, along $\Delta$ and in the face, and taking
$m=\hbar^2({\mathrm{d}}^2E/{\mathrm{d}} k^2)^{-1}$.  Rosei's masses come from exactly
the same source (``the optical masses and the gaps\ldots were taken
from Ref.~4'', his Ref.~4 being C\&S), so the two readings of the same
figure ought to agree.  They do not: my extraction gives
$m_{u\parallel}=0.40$, $m_{l\parallel}=0.15$ --- the \emph{opposite}
$\parallel$ ordering to Rosei's, hence ${\bar{B}}_X>0$.  On that branch the
sub-gap picture changes: the integration window becomes \emph{closed}
--- finite at both ends by geometry, as at $L$ --- instead of Rosei's
bottomless, occupation-cut window; there is no geometric sub-gap tail,
and the onset is a $\sqrt{{\hbar\omega}-{\mathcal{E}_g^{X}}}$ rise sitting exactly at the gap
(any observed tail below $1.94$ eV would then be broadening only, and
the $-20{k_B} T$ floor would overshoot the true, geometric one).

Two facts keep the fork contained:
\begin{itemize}
\item Above the gap the branches are \emph{indistinguishable}: same line shape, same singular edge.  The $\parallel$ masses enter only through ${\mathcal{F}}_X$, and the fit lumps ${\mathcal{F}}_X$ into $S_X={\mathcal{F}}_X|P_X|^2$ --- so no fit to ${\varepsilon}_2$ data can tell the branches apart; only sub-gap behaviour or the extracted $|P_X|^2$ could.  (The above-gap window slope $A_u/{\bar{A}}_X$ --- $0.62$ on Rosei's assignment, $0.38$ on mine --- is set by the $\perp$ masses, a separate assignment.)
\item Either ordering makes the $d$ band the \emph{lighter} one along $\Delta$, which goes against the flat-$d$ prior --- so the likeliest culprit is a transcription swap in my own mass table, not new physics.
\end{itemize}

\noindent\textbf{Bottom line: until C\&S Fig.~5 is re-measured, this
document follows Rosei's branch (${\bar{B}}_X<0$) everywhere, and quotes the
branch wherever it matters.}

% =====================================================================
\section{Assembled ${\varepsilon}_2$}
\label{sec:eps2}
% =====================================================================

Insert the EDJDOS ($4'$) into the FGR form \eqref{eq:star}.  The
bridge is the definition of ${D_{l\to u}}$ itself, Eq.~\eqref{eq:sort}: the
$(2\pi)^3$ cancels, and with spin (the explicit 2) and the
half-neighborhood counts $N_X=6$, $N_L=8$,
${\varepsilon}_2$
which tidies to Rosei's second headline equation, with no leftover
constant:
$\to$
<!--\label{eq:eps9}-->
(Rosei leaves the $N_i$ implicit, so the fitted $|P_i|^2$ really means
$N_i|P_i|^2$ up to normalization --- harmless for the quoted ratio).
That clean landing needs three things at once: $|P|^2$ is the squared
gradient element \eqref{eq:Pdef}; ${D_{l\to u}}$ is ($4'$), not the printed
(4); spin and $N_i$ are kept where shown.  The quantities the data
actually fix are the strengths
${k_{\parallel}}$
<!--\label{eq:S10}-->
with $|P_X/P_L|^2=0.370$ from the fit to Johnson--Christy; under the
corrected ($4'$) the fitted combination maps onto the printed $S_i$
through exactly the constant $2{\mathcal{F}}_i/\hbar^2$ of
Sec.~\ref{sec:kernel} --- every published number keeps its meaning.

Fully substituted --- the explicit integral a computer evaluates ---
($4'$) into (9), the $\hbar^4$'s cancelling:
$\mathbf{k}$
<!--\label{eq:eps11}-->
with ${k_{\parallel}}^{X}$ from ($6'$), ${k_{\parallel}}^{L}$ from its $L$ twin (same
bracket, every slot sign-reversed), and windows:

\begin{center}
\small
\begin{tabular}{lccl}
\toprule
 & $E_{\min}$ & $E_{\max}$ & edge behaviour \\
\midrule
$X$, ${\hbar\omega}\ge{\mathcal{E}_g^{X}}$ & $-20{k_B} T$ (occupation) & ${\hbar\omega_{X_6^-}}+\frac{A_u}{{\bar{A}}_X}({\hbar\omega}-{\mathcal{E}_g^{X}})$ & $1/\sqrt{\;}$ at $E_{\max}$ \\
$X$, ${\hbar\omega}<{\mathcal{E}_g^{X}}$ & $-20{k_B} T$ (occupation) & sub-gap line of ($8'$) & finite everywhere \\
$L$, ${\hbar\omega}\ge{\mathcal{E}_g^{L}}$ & ${\hbar\omega_{L_6^-}}+\frac{A_u}{{\bar{A}}_L}({\hbar\omega}-{\mathcal{E}_g^{L}})$ & ${\hbar\omega_{L_6^-}}+\frac{B_u}{{\bar{B}}_L}({\hbar\omega}-{\mathcal{E}_g^{L}})$ & $1/\sqrt{\;}$ at $E_{\min}$ \\
$L$, ${\hbar\omega}<{\mathcal{E}_g^{L}}$ & --- & --- & no states \\
\bottomrule
\end{tabular}\end{center}

\medskip
\noindent\textbf{Domain of validity --- where is $k_{\max}$?}
The parabolic bands \eqref{eq:bands} and the constant $|P_i|^2$ hold
only near the critical points; Rosei sizes the trusted region as
${\sim}(\pi/10a)^3$ (matrix-element constancy, his Sec.~IV), while his
Figs.~3 and 5 draw the CEDS out to $k\sim2\times\pi/4a$.  Yet none of
the windows above carries an explicit $k_{\max}$ --- they encode CEDS
geometry and occupation only.  This is faithful to the paper: Rosei
polices locality through the \emph{photon energy} instead, fitting
only to 2.7 eV --- ``only limited regions in $K$ space around $X$ and
around $L$ are involved \dots\ when the fitting is extended a few
tenths of an eV above the absorption edge'', while ``the danger of
counting transitions twice and of a breakdown of our approximations''
(both channels are band-5\,$\to$\,band-6) ``may therefore arise at
higher photon energies''.  The numbers behind that policy, with the
parameters below ($a_{\rm Au}=4.08$~\AA, so
$\pi/4a=0.19$~\AA$^{-1}$):
\begin{itemize}
\item \emph{$L$ is safely local}: the CEDS ellipsoid lies entirely inside $|\mathbf k|\le\pi/4a$ for ${\hbar\omega}-{\mathcal{E}_g^{L}}\le{\bar{A}}_L(\pi/4a)^2\simeq0.79$ eV, i.e.\ up to ${\hbar\omega}\simeq3.2$ eV --- the whole fit range (inside $\pi/10a$ only for ${\hbar\omega}\lesssim2.6$ eV).
\item \emph{$X$ is marginal by construction}: the occupied edge of the window ($E\simeq E_F$, where the Fermi step sits) lies at $|\mathbf k|\simeq(0.8\text{--}1.5)\times\pi/4a$ for ${\hbar\omega}=2.0$--$2.7$ eV --- several times the strict constant-$|P|^2$ radius, and exactly the scale of Rosei's own figures.  The $-20{k_B} T$ floor is an \emph{occupation} cutoff, not a locality one --- it maps to $|\mathbf k|$ up to ${\sim}1.9\times\pi/4a$ --- but in equilibrium everything beyond the Fermi step carries $[1-{f}]\approx0$, so the occupation acts as the de facto $k$-cutoff.  Nothing in the model \emph{checks} this.
\end{itemize}
Operationally: trust the model where Rosei did, ${\hbar\omega}\lesssim2.7$ eV;
above that, locality and single-counting fail together.  Out of
equilibrium the occupation excuse at $X$ weakens --- see the locality
guard in Sec.~\ref{sec:numerics}.

\medskip
\noindent\textbf{Parameters} (levels: Rosei's fit; masses: C\&S, the
$X$ row on Rosei's branch, subject to Sec.~\ref{sec:fork}; ${\mathcal{F}}$ in
units of the free-electron mass $m$; ${\mathcal{D}}$ in units $(\hbar^2/2m)^2$):

\begin{center}
\small
\begin{tabular}{lccccccc}
\toprule
 & ${\mathcal{E}_g}$ (eV) & upper level (eV) & lower level (eV) & $m_{u\perp}$ & $m_{u\parallel}$ & $m_{l\perp}$ & $m_{l\parallel}$ \\
\midrule
$X$ & 1.94 & ${\hbar\omega_{X_6^-}}=0.17$ & ${\hbar\omega_{X_7^+}}=1.77$ & 0.19 & 0.15 & 0.31 & 0.40 \\
$L$ & 2.45 & ${\hbar\omega_{L_6^-}}=0.89$ & ${\hbar\omega_{L^+}}=1.56$ & 0.24 & 0.12 & 0.70 & 1.03 \\
\bottomrule
\end{tabular}
\medskip
\begin{tabular}{lcccc}
\toprule
 & ${\mathcal{D}}$ & ${\mathcal{F}}/m$ & window slopes & singular edge \\
\midrule
$X$ & $+34.7$ & 0.170 & $E_{\max}$ (above gap): 0.62 & upper (${k_{\parallel}}\to0$) \\
$L$ & $-7.86$ & 0.357 & $E_{\min}$: 0.745,\ $E_{\max}$: 0.896 & lower (${k_{\parallel}}\to0$) \\
\bottomrule
\end{tabular}\end{center}

\medskip
\noindent\textbf{Scorecard --- the ten equations, one by one:}

\begin{center}
\small
\begin{tabular}{llll}
\toprule
Rosei & here & status & issue as printed \\
\midrule
(1) & (1) & recreated & --- \\
(2) & (2) & recreated & --- \\
(3) & (3) & recreated & --- \\
(4) & ($4'$) & \textbf{suggested fix} & prefactor $(8\pi^2\hbar^2)^{-1}{\mathcal{F}}\to(4\pi^2\hbar^4)^{-1}{\mathcal{F}}^2$; off by const.\ $2{\mathcal{F}}/\hbar^2$ \\
(5) & (5) & recreated & --- \\
(6) & ($6'$) & \textbf{suggested fix} & dimensionally impossible as printed \\
(7) & (7) & recreated & --- \\
(8) & ($8'$) & \textbf{suggested fix} & $m_{u\parallel}\leftrightarrow m_{l\parallel}$ subscripts swapped \\
(9) & (9) & recreated & $N_i$ made explicit \\
(10) & (10) & recreated & absorbs the (4) constant; see App.~\ref{app:dict} \\
\bottomrule
\end{tabular}\end{center}

% ---------------------------------------------------------------------
\subsection{Discrepancies \& suggested corrections}\label{sec:grw}
% ---------------------------------------------------------------------

\noindent\textbf{At issue: the printed Eqs.~(4), (6) and (8) of the
1975 paper.}  None of the three comes out of the rederivation as
printed --- and none of the three failures touches a published number.
Each failure stated simply, then fixed.

\medskip
\textbf{Eq.\ (4).}  Printed:
${D_{l\to u}}=(8\pi^2\hbar^2)^{-1}\,{\mathcal{F}}_{l\to u}\,{k_{\parallel}}^{-1}$.
\emph{Why it fails:} the definition of ${D_{l\to u}}$ (Eq.~\eqref{eq:sort})
leaves nothing to choose --- the reduction of
Sec.~\ref{sec:kernel} gives $(16\pi^2|{\mathcal{D}}|)^{-1}{k_{\parallel}}^{-1}
={\mathcal{F}}^2/(4\pi^2\hbar^4{k_{\parallel}})$, Eq.~($4'$), and with the gradient
convention for $|P|^2$ only ($4'$) makes \{(4),(9)\} dimensionally
consistent and lands (9) with no leftover constant.  The printed form
is off by
$\mathbf{A}$
a pure constant per critical point --- it cannot bend a line shape,
and it is absorbed once and for all into the fitted $S={\mathcal{F}}|P|^2$ of
Eq.~(10).  That is why the slip is easy to miss and costs nothing
physically; but the assembly only closes into (9) if ($4'$) is used.

\medskip
\textbf{Eq.\ (6).}  Printed:
$\nabla\!\cdot\!\mathbf{A}=0$
\emph{Why it fails:} a wavevector is set equal to the square root of a
sum of \emph{energies}, and inside the bracket a bare energy is added
to $(\hbar^2/2m)\times$energy terms --- mismatched units throughout;
the coefficients read $\hbar^2/2m$ where $2m/\hbar^2$ is needed for
the bracket to come out as length$^{-2}$.  Three small fixes turn it
into ($6'$): square the left-hand side; invert the coefficients and
pull out the overall $2{\mathcal{F}}^2/\hbar^2$; restore the missing
$1/m_{u\perp}$ on the first term.  The three numerator slots of
$\hat{\mathbf p}\!\cdot\!\mathbf A
=\mathbf A\!\cdot\!\hat{\mathbf p}$
then line up with the print term by term --- confirmation that ($6'$)
is what was meant: a typo, but a disabling one, since the printed
formula is unusable as written.

\medskip
\textbf{Eq.\ (8).}  Printed:
$E_{\max}={\hbar\omega_{X_6^-}}+({\mathcal{E}_g^{X}}-{\hbar\omega})\,m_{u\parallel}/(m_{l\parallel}-m_{u\parallel})$.
Solving the model for the ${k_{\perp}}=0$ edge gives
$\tfrac12$
i.e.\ the printed equation has $m_{u\parallel}$ and $m_{l\parallel}$
interchanged --- one consistent swap, specific to (8): the subscripts
of (6) are \emph{not} swapped, so it is not the same slip repeated.
\emph{The tell:} on Rosei's own branch $m_{u\parallel}<m_{l\parallel}$,
the printed denominator is positive and sends $E_{\max}$ \emph{above}
${\hbar\omega_{X_6^-}}$ as ${\hbar\omega}$ drops below the gap --- but the sub-gap CEDS has no
states up there.  ($8'$) puts the sub-gap edge sensibly \emph{below}
${\hbar\omega_{X_6^-}}$, which is what makes the tail (and the $1.83$ eV step onset)
work.  On the alternative branch of Sec.~\ref{sec:fork}
(${\bar{B}}_X>0$), the same algebraic edge is instead the \emph{lower}
window limit.

% ---------------------------------------------------------------------
\subsection{The Drude term: completing ${\varepsilon}_2$}\label{sec:drude}
% ---------------------------------------------------------------------

Everything above is interband.  The measured ${\varepsilon}_2$ of gold also
contains the free-carrier (intraband) response, which the Rosei model
deliberately leaves out --- so to compare with data, or to fit, it has
to be added back.  The Drude dielectric function
${\varepsilon}(\omega)={\varepsilon}_\infty-\omega_p^2/[\omega(\omega+i\gamma)]$ gives
$A^2$
<!--\label{eq:drude}-->
and in the optical window the high-frequency limit is all that is
needed: for Au, $\hbar\omega_p\simeq9.0$ eV and
$\hbar\gamma\simeq0.07$ eV, so $\omega\gg\gamma$ holds everywhere
above $\sim\!1$ eV (the fit code implements exactly this
$A_{\rm D}/\omega^3$ form, $A_{\rm D}=\omega_p^2\gamma\simeq5.7$
eV$^3$).  The total is a plain sum of channels:
$V$
<!--\label{eq:total}-->
The Drude part falls as $\omega^{-3}$ and carries no interband
structure: below the $X$ onset it \emph{is} ${\varepsilon}_2$, through the edge
region it is a smooth declining background, and the temperature
dependence of the edge sits entirely in the ${\mathcal{J}}_i$.

% =====================================================================
\section{Photonic density of states and the emission counterpart}
\label{sec:photonic}
% =====================================================================

\subsection{Mode counting}

Periodic quantization in volume $V$: two transverse polarizations,
$\omega=cq$:
$\hat H_{\rm int}$
<!--\label{eq:rhophot}-->
modes per unit volume per unit angular frequency.  In a transparent
host of index $n$, $c\to c/n$; in a structured environment
$\rho\to\rho(\mathbf r,\omega)$, the \emph{local} density of states
(Novotny \& Hecht) --- precisely the photonic factor of the PL
factorization used in this thesis.

\subsection{Spontaneous emission rate of one pair}

FGR with the quantized field \eqref{eq:Aquant} at $n_q=0$, summed over
modes, using
$\sum_\lambda\int{\mathrm{d}}\Omega_q|\hat{\bm{\varepsilon}}_\lambda\!\cdot\!{\mathbf{p}_{ul}}|^2
=\tfrac{8\pi}{3}|{\mathbf{p}_{ul}}|^2$ and $|{\mathbf{p}_{ul}}|^2=\hbar^2|P|^2$:
$+1$
<!--\label{eq:Wsp}-->
--- the vacuum Einstein-$A$ rate in the
$\mathbf A\!\cdot\!\mathbf p$ form.  The factorization into a matter
part times $\rho(\omega)$ is exact and is the microscopic origin of
the LDOS $\times$ electronic split assumed in Chapter 2.

Why $1/\omega$ in the matter part here, when the absorption prefactor
of \eqref{eq:eps2exact} carries $1/\omega^2$?  Both rates contain the
identical vertex $(e/mc)^2|A|^2|M_{ul}|^2$; the different powers come
from what $|A|^2$ is normalized against.  In absorption, ${\varepsilon}_2$ is
defined per unit \emph{classical field intensity}: expressing
$A_0=cE_0/\omega$ through the measured field costs two inverse powers
of $\omega$ (the $-2$ of Eq.~\eqref{eq:wcount}).  In spontaneous
emission there is no external field to normalize by; the reference
object is \emph{one photon per mode}, whose squared amplitude is fixed
by the field quantization \eqref{eq:Aquant} at $2\pi\hbar c^2/\omega
V$ --- a \emph{single} inverse power of $\omega$, the zero-point
amplitude of the mode.  The remaining $\omega$-dependence of
\eqref{eq:Wsp} is carried openly by the mode count
$\rho(\omega)\propto\omega^2$.  In short: same vertex, different
yardstick --- absorption is measured against $E_0^2$ (two powers),
emission against the vacuum fluctuation of one mode (one power).

\subsection{The interband PL spectrum and the Kirchhoff check}

The emission weight is not a difference but a \emph{product}: the
upper level must be filled and the lower empty,
${f}(E)[1-{f}(E-{\hbar\omega})]$.  Summing \eqref{eq:Wsp} over pairs (spin 2,
$N$ half-neighborhoods, same ${D_{l\to u}}$ and windows as in absorption):
$n_q=0$
<!--\label{eq:PL}-->
photons per unit time, volume and photon energy.  In equilibrium the
Fermi identity $[1-{f}]/{f}={\mathrm{e}}^{(E-\mu)/{k_B} T}$, applied at both
levels, collapses the product (detailed balance; only the
\emph{transition} energy survives):
$\mathbf A$
<!--\label{eq:BEcollapse}-->
and since $[{f}(E-{\hbar\omega})-{f}(E)]$ is Rosei's $[1-{f}(E)]$ in the same
deep-$d$-band limit used throughout (Sec.~\ref{sec:occupation}), every
electronic factor of \eqref{eq:PL} maps onto (11) and the constants
collapse \emph{exactly}:
$\psi_{n\mathbf k}={\mathrm{e}}^{i\mathbf k\cdot\mathbf r}
u_{n\mathbf k}(\mathbf r)/\sqrt V$
<!--\label{eq:vRS}-->
--- the van Roosbroeck--Shockley relation; per thermal photon,
emission over absorption is ${\mathrm{e}}^{-{\hbar\omega}/{k_B} T}$: Kirchhoff.  That
\eqref{eq:eps9} and \eqref{eq:PL}, derived independently, satisfy
\eqref{eq:vRS} with no leftover constant is the strongest global check
that every prefactor in this document is right.

Out of equilibrium the collapse \eqref{eq:BEcollapse} fails --- it
used one shared $(\mu,T)$ at both levels --- and the ratio
$n_{\rm eff}(E,{\hbar\omega})={f}(E)[1-{f}(E-{\hbar\omega})]/[{f}(E-{\hbar\omega})-{f}(E)]$
retains genuine $E$-dependence: the emission line shape parts company
with ${\varepsilon}_2$.  That deviation \emph{is} the non-equilibrium PL signal
this thesis is after; reproducing $n_{\rm eff}\to n_B$ under
${f}_{\rm FD}$ is a unit test, \emph{taking} it is throwing the
signal away.  In a structured environment,
$\rho(\omega)\to\rho(\mathbf r,\omega)$ in \eqref{eq:PL}.

% =====================================================================
\section{Numerics: evaluate in ${k_{\parallel}}$, not in $E$}\label{sec:numerics}
% =====================================================================

Groundwork for the implementation.  The $1/\sqrt{\;}$ edges of (11)
would force adaptive quadrature or endpoint-weighted rules --- but the
divergence is a pure coordinate artifact, manufactured in
Sec.~\ref{sec:kernel} when the smooth $\mathbf k$-space geometry was
re-parametrized by final energy.  One exact change of variable inside
the same 1D integral undoes it.  Take $k\equiv{k_{\parallel}}$ itself.  (Where
did ${k_{\perp}}$ go?  Nowhere --- on the transition surface it is not an
independent variable.  With ${\hbar\omega}$ fixed, the constraint
${\mathcal{E}_{lu}}={\hbar\omega}$ of \eqref{eq:Omega} slaves the perpendicular
coordinate, ${k_{\perp}}^{2}=({\hbar\omega}-{\mathcal{E}_g}-{\bar{B}}{k_{\parallel}}^{2})/{\bar{A}}$, exactly as $u$
was expressed through $v$ in \eqref{eq:uv}; substituting it into
$E=\hbar\omega_u(\mathbf k)$ of \eqref{eq:bands} leaves $E$ a function
of ${k_{\parallel}}$ alone, Eq.~\eqref{eq:Eofk} below --- ${k_{\perp}}$ is
eliminated by the $\delta$-constraint, not set to zero.)  Then $E$ is
affine in $k^2$, so ${\mathrm{d}} E=\mp(\mu_{i\perp}\hbar^2/{\mathcal{F}}_i^2)\,k\,{\mathrm{d}} k$
while ${D_{l\to u}}\propto1/k$ --- the factors of $k$ cancel identically,
${\mathrm{e}}^{\pm i\mathbf q\cdot\mathbf r}\,\hat{\bm{\varepsilon}}\!\cdot\!\hat{\mathbf p}$
i.e.\ ${k_{\parallel}}$ is (up to constants) the cumulative count of the
EDJDOS: integrating in ${k_{\parallel}}$ \emph{is} integrating in $E$ with the
known weight ${D_{l\to u}}$ absorbed exactly into the node placement.
Eq.~(7) flattens to
$\mathbf k'=\mathbf k\pm\mathbf q$
<!--\label{eq:k12}-->
$q\sim\omega/c\sim10^{-3}\,$
<!--\label{eq:Eofk}-->
valid on both sides of the $X$ gap --- all the branch logic lives in
the limits.  These are the $u=0$/$v=0$ endpoints of
Sec.~\ref{sec:XL}, written out since an implementation touches them
first.  At $X$ (Rosei's branch):
$^{-1}$
with $E_{\rm floor}=-20{k_B} T$ in equilibrium; for a pumped
distribution take $E_{\rm floor}=-(\hbar\omega_{\rm pump}+20{k_B} T)$
--- generosity is cheap, since below the true floor the integrand is
${\approx}\,0$.  At $L$ both limits are geometric:
$\mathbf k'=\mathbf k$

No singularity survives: the integrand of (12) is a bounded, smooth
occupation profile whose only feature is the Fermi step of width
$\sim{k_B} T$.  \textbf{Recipe} (composite Simpson on a uniform grid is
entirely adequate):
\begin{enumerate}
\item $k=\mathrm{linspace}(k_1({\hbar\omega}),\,k_2({\hbar\omega}),\,N)$, $N$ odd;
\item integrand $1-{f}(E_i(k))$ --- Fermi--Dirac in equilibrium, interpolation of a tabulated non-equilibrium distribution otherwise (use a \emph{monotone} interpolant, e.g.\ PCHIP, to avoid overshoot at the distribution's kinks);
\item \texttt{simpson(y, x=k)}, times $\mu_{i\perp}/4\pi^2\hbar^2$; then into (9).
\end{enumerate}
Measured accuracy (equilibrium, ${\hbar\omega}=2.4$ eV, against a converged
reference): $X$ reaches $5\times10^{-10}$ relative error at $N=101$
and machine precision by $N=201$; the $L$ integrand in equilibrium is
constant across the window (Remark~\ref{rem:neck}), and even a hot
2000 K distribution gives $3\times10^{-14}$ at $N=101$.  By contrast,
uniform-grid Simpson in $E$ on a singular edge stalls near $10^{-2}$
even at $N=6401$ ($O(\sqrt h)$ decay).  Energy space \emph{done
properly} --- the substitution $t=\sqrt{E_{\rm edge}-E}$, or
Gauss--Jacobi with the $1/\sqrt{\;}$ weight --- is bit-identical to
(12), because $t\propto{k_{\parallel}}$ exactly: the geometry had already named
the correct quadrature variable.  Hence the division of labor:
\emph{derive and reason in $E$; evaluate in ${k_{\parallel}}$.}

\vspace{1em}
\textbf{Practical points.}
\begin{itemize}
\item \emph{Resolve the step}: the only feature sits at $k_{\rm step}=k(E{=}0)$, trivially located since $E_i(k)$ is affine in $k^2$; at low $T$ either raise $N$ or split the interval there and Simpson each panel.
\item \emph{Respect the kinks} of a pumped distribution: slope discontinuities at $E_F\pm\hbar\omega_{\rm pump}$ are \emph{distribution} features no variable change removes --- put panel boundaries at their images $k(E_F\pm\hbar\omega_{\rm pump})$.
\item \emph{Vectorize over ${\hbar\omega}$}: build $k=k_1[:,\!None]+(k_2-k_1)[:,\!None]\,\xi[None,:]$ with $\xi$ uniform on $[0,1]$ and Simpson along the last axis --- the whole spectrum in one call.  Gauss--Legendre with 48--64 nodes does the same job with fewer points if node count matters.
\item \emph{Cache the geometry}: $k_{1,2}$, $E_i(k)$ and the grid depend on $({\hbar\omega})$ only --- not on $T$ or on the distribution --- so precompute them once and sweep temperatures/pump powers cheaply (this is what makes fitting fast).
\item \emph{Locality guard}: nothing in (12) knows where the parabolic neighbourhood ends --- compare $k_2({\hbar\omega})$ with $\pi/4a\simeq0.19$~\AA$^{-1}$ (Sec.~\ref{sec:eps2}).  In equilibrium the $X$ integrand dies of occupation before locality does; a pumped distribution re-weights exactly those deep-window nodes, and the generous floor pushes $k_2$ to ${\sim}2\times\pi/4a$.  Report the fraction of (12) accrued at $k>\pi/4a$ and distrust spectra where it is not small.
\item \emph{Error control}: recompute at $2N$ and compare (one doubling).
\item \emph{Emission}: the PL integrand \eqref{eq:PL} swaps $[1-{f}]\to {f}(E)[1-{f}(E-{\hbar\omega})]$ in the same quadrature --- geometry, windows and nodes unchanged.
\end{itemize}

\textbf{Built-in checks.}
(i) \emph{Empty-window arithmetic}: wherever the occupation factor is
$\equiv1$ across the window,
${\mathcal{J}}_i=\frac{\mu_{i\perp}}{4\pi^2\hbar^2}(k_2-k_1)$ exactly; at $L$
with $T\to0$ this is the textbook $M_0$ onset
${\mathcal{J}}_L=\frac{\mu_{L\perp}}{4\pi^2\hbar^2}
\sqrt{({\hbar\omega}-{\mathcal{E}_g^{L}})/{\bar{B}}_L}$ in closed form.
(ii) \emph{Two-forms agreement}: (11) by adaptive quadrature (endpoint
declared) and (12) by Simpson must agree to quadrature accuracy at
every $({\hbar\omega},T)$ --- a one-line unit test.
(iii) \emph{Kirchhoff}: with ${f}={f}_{\rm FD}$, confirm
$n_{\rm eff}(E,{\hbar\omega})\to n_B({\hbar\omega})$ independent of $E$
(Sec.~\ref{sec:photonic}).

% =====================================================================
\appendix
% =====================================================================
\section{Conventions}\label{sec:conv}

Fixed once here, used silently throughout the document.

\begin{center}
\begin{tabular}{ll}
\toprule
Units & Gaussian throughout (Rosei's system); SI conversions noted where useful. \\
\midrule
Energy zero & Fermi level, $E_F=0$; all electron energies $E$ from $E_F$. \\
Bands & $u$ = upper (conduction), $l$ = lower ($d$); Rosei's subscripts. \\
Levels \& masses & ${\hbar\omega_{X_6^-}},{\hbar\omega_{X_7^+}},{\hbar\omega_{L_6^-}},{\hbar\omega_{L^+}}$ and $m_{u\perp},m_{u\parallel}, m_{l\perp},m_{l\parallel}$ are all \\
 & \emph{positive numbers}; curvature directions and the side of $E_F$ are carried \\
 & only by explicit $\pm$ signs, never hidden in a symbol. \\
Band coefficients & $A_i=\hbar^2/2m_{i\perp}$, $B_i=\hbar^2/2m_{i\parallel}$, $i=u,l$; all $A_i,B_i>0$. \\
Interband energy & ${\mathcal{E}_{lu}}(\mathbf k)\equiv\hbar\omega_u-\hbar\omega_l$; Rosei's $\Omega_{lu}={\mathcal{E}_{lu}}-{\hbar\omega}$, the CEDS being $\Omega_{lu}=0$. \\
Local axes & ${k_{\parallel}}$ along $\Gamma X$ (resp.\ $\Gamma L$), ${k_{\perp}}$ radial in the perpendicular plane; \\
 & $\mathbf{k}$ measured from the critical point. \\
Spin & carried explicitly as a factor 2, never absorbed silently. \\
Counting & $N$ counts equivalent \emph{half}-neighborhoods ($N_X=6$ faces, $N_L=8$), \\
 & paired with the half-space ${D_{l\to u}}$ below --- or a factor 2 goes missing. \\
Occupation & Rosei's $[1-{f}(E,T)]$, ${f}(E)=[{\mathrm{e}}^{E/{k_B} T}+1]^{-1}$; see Sec.~\ref{sec:occupation}. \\
Matrix element & $\vert{}P\vert{}^2\equiv\vert{}\langle u\vert{}\nabla\vert{}l\rangle\vert{}^2$ (dimension length$^{-2}$), Rosei's convention; \\
 & momentum form $\vert{}{\mathbf{p}_{ul}}\vert{}^2=\hbar^2\vert{}P\vert{}^2$. \\
Photon & energy ${\hbar\omega}$, wavevector $\mathbf{q}$, polarization $\hat{\bm{{\varepsilon}}}_\lambda$. \\
\bottomrule
\end{tabular}\end{center}

\section{Dictionary of conventions across the repository}\label{app:dict}

\begin{center}
\footnotesize
\setlength{\tabcolsep}{4pt}
\begin{tabular}{lllll}
\toprule
Document & upper band & lower band & ${\bar{B}}$ & status \\
\midrule
Rosei 1975; this doc ($X$) & $+A_u{k_{\perp}}^2-B_u{k_{\parallel}}^2$ & $-A_l{k_{\perp}}^2-B_l{k_{\parallel}}^2$ & $B_l-B_u$ & canonical \\
\emph{Equivalence} note ($X$) & $+A_u{k_{\perp}}^2+B_u{k_{\parallel}}^2$ & $-A_l{k_{\perp}}^2+B_l{k_{\parallel}}^2$ & $B_u-B_l$ & same ${\mathcal{D}}$ (sign-robust) \\
A8 / A8.L & $+A_c u-B_c v$ & $-A_v u-B_v v$ & $B_v-B_c$ & A8 ok$^{(a)}$; A8.L wrong at $L$$^{(b)}$ \\
\texttt{interband\_derivation\_v2} & mixed & mixed & --- & internally inconsistent$^{(c)}$ \\
\bottomrule
\end{tabular}\end{center}
{\footnotesize
$^{(a)}$ A8's $\mathcal{E}_{\min}$ denominator $(B_v{+}B_c)$ is the
$s{=}{+}1$ formula; for its own $X$ conventions it should be
${\bar{B}}=B_v{-}B_c$.\quad
$^{(b)}$ A8.L models the $L$ conduction band as a saddle ($-B_c$);
C\&S give $m_{u\parallel}^L=+0.12$: it is a minimum, $s=+1$.\quad
$^{(c)}$ v2 asserts ${\bar{B}}>0$ \emph{and} a geometric sub-gap tail
(requires ${\bar{B}}<0$), and labels $L$ as M$_1$.}

\medskip
\noindent\textbf{Normalizations.}  Earlier repo documents use the
full-neighborhood, no-$(2\pi)^3$ kernel
$K(E,{\hbar\omega})=\int{\mathrm{d}}^3k\,\delta\delta$ and the half-line measure
$\mathcal J_{\rm eq}$ of the \emph{Equivalence} note; with ${D_{l\to u}}$ from
($4'$),
${\varepsilon}_2$
Matrix elements: Rosei's $|P|^2$ is the squared \emph{gradient}
element, related to the SI momentum element by
$|{\mathbf{p}_{ul}}|^2=\hbar^2|P|^2$.  With these two conventions --- ($4'$) and
gradient $|P|^2$ --- Rosei's Eq.~(9) is exact as printed; no residual
constant needs to be hidden in $|P|^2$, and only the strengths
$S={\mathcal{F}}|P|^2$ of Eq.~(10) are fixed by the data (the printed-(4)
constant reshuffles ${D_{l\to u}}$, ${\mathcal{F}}$, $|P|^2$ among themselves without
touching $S$).

\section{Dimensional audit (Gaussian)}\label{app:dim}

$[\,e^2\,]={\rm erg\,cm}$, $[\,|P|^2\,]={\rm cm^{-2}}$,
$[\,A_i\,]=[\,B_i\,]={\rm erg\,cm^2}$, $[\,{\mathcal{D}}\,]={\rm erg^2cm^4}$,
$[\,{\mathcal{F}}\,]={\rm g}$, $[\,{D_{l\to u}}\,]={\rm erg^{-2}cm^{-3}}$,
$[\,{\mathcal{J}}\,]={\rm erg^{-1}cm^{-3}}$.
\begin{itemize}
\item ($4'$): ${\mathcal{F}}^2/\hbar^4{k_{\parallel}}\sim {\rm g^2}/({\rm erg^4s^4\,cm^{-1}})={\rm erg^{-2}cm^{-3}}$ --- states per volume per (energy)$^2$, as a double-$\delta$ measure must be. \checkmark
\item (9): $\dfrac{{\rm erg\,cm}\cdot{\rm erg^4s^4}} {{\rm g^2}\,{\rm erg^2}}\cdot{\rm cm^{-2}}\cdot {\rm erg^{-1}cm^{-3}}=1$ (using ${\rm erg\,s^2=g\,cm^2}$). \checkmark
\item (12): $\mu_\perp/\hbar^2\times[k]\sim {\rm g\,cm^{-1}}/{\rm erg^2s^2}={\rm erg^{-1}cm^{-3}}=[{\mathcal{J}}]$. \checkmark
\item \eqref{eq:rhophot}: $[\rho]={\rm s\,cm^{-3}}$. \checkmark
\item \eqref{eq:Wsp}: $\dfrac{{\rm erg\,cm}\cdot{\rm s^{-1}}\cdot {\rm erg\,s}\cdot{\rm cm^{-2}}}{{\rm g^2\,cm^3s^{-3}}} ={\rm erg^2\,cm^{-1}}/({\rm g^2cm^3s^{-3}}\,{\rm s}) ={\rm s^{-1}}$. \checkmark
\item \eqref{eq:PL}: $[\rho]\cdot[\text{matter}]\cdot[{\mathcal{J}}] ={\rm s\,cm^{-3}}\cdot{\rm cm^3s^{-2}}\cdot{\rm erg^{-1}cm^{-3}} ={\rm erg^{-1}cm^{-3}s^{-1}}$ --- photons per time, volume and photon energy. \checkmark
\item \eqref{eq:vRS}: $[\omega\rho/\hbar] ={\rm erg^{-1}cm^{-3}s^{-1}}$, times dimensionless ${\varepsilon}_2n_B$. \checkmark
\end{itemize}

\section*{References}
\begin{itemize}
\item M.~Guerrisi, R.~Rosei, P.~Winsemius, \emph{Splitting of the interband absorption edge in Au}, Phys.\ Rev.\ B \textbf{12}, 557 (1975) --- the paper recreated here ($X$ derived, $L$ delegated); source of the $X$-point dispersion relations, Eqs.~(1)--(2).
\item R.~Rosei, \emph{Temperature modulation of the optical transitions involving the Fermi surface in Ag: theory}, Phys.\ Rev.\ B \textbf{10}, 474 (1974) --- the $L$ machinery (and the EDJDOS recipe) the 1975 paper borrows; source of the $L$-point dispersion relations.
\item N.~E.~Christensen, B.~O.~Seraphin, \emph{Relativistic band calculation and the optical properties of gold}, Phys.\ Rev.\ B \textbf{4}, 3321 (1971) --- source of the levels and masses (Rosei's Ref.~4).
\item P.~B.~Johnson, R.~W.~Christy, \emph{Optical constants of the noble metals}, Phys.\ Rev.\ B \textbf{6}, 4370 (1972) --- the data the fit targets.
\item Y.~Sivan, Y.~Dubi, non-equilibrium electron distributions and metal photoluminescence (2021), and refs.\ therein.
\item L.~Novotny, B.~Hecht, \emph{Principles of Nano-Optics}, 2nd ed., Cambridge (2012) --- LDOS.
\end{itemize}

\end{document}
